\chapter{重力与质量：测量力}\label{ch:g9:weight-vs-mass}

五年前，在纸上的月球上，你把“重”一分为二：
质量描述物质，重力描述拉力。这一划分很明智；
如今它要变成计算。力有了自己的单位——
以苹果落地故事中的牛顿命名——重力有了公式，
童年时的橡皮筋秤，
以专业的面貌重新出现。

\section{牛顿}

\begin{definition}[牛顿]\label{def:g9:weight-vs-mass:newton}
力——推、拉、重力、万有引力的作用——都以
\emph{牛顿}\index{牛顿}（\unit{N}）为单位计量，以纪念
万有引力定律的提出者。便于估计的参照：
一牛顿约是一个小苹果的重量（历史在此莞尔）；
你的书包约重 \qty{50}{N}；一次有力的握手
约用 \qty{100}{N}；上一章的月球租约
达到 $10^{20}$。
\end{definition}

\begin{definition}[测力计]\label{def:g9:weight-vs-mass:dynamometer}
\emph{测力计}\index{测力计}是一种力的测量仪：一个
标定过的弹簧装在外壳里，伸长量与所受的力成正比，
指针直接以牛顿为读数。它是童年橡皮筋秤的长大版——
与此前的电压表、电流表一样，它让又一块
“感觉很大”的印象
变成了可靠的数字。
\end{definition}

\begin{method}[测量力]\label{met:g9:weight-vs-mass:measure}
\begin{enumerate}
  \item 选择量程与任务相配的测力计——
        量程 \qty{2}{N} 的用于铅笔盒，
        \qty{100}{N} 的则用于书包；
  \item 在使用时的位置上调零（若负载是悬挂的，
        就应在悬挂状态下调零——弹簧自身的重量
        会使零点发生偏移）；
  \item 平稳地施加力，沿仪器的
        轴线方向施力，并在与视线齐平处读数；
  \item 连同单位记录读数——由于力有方向，
        必要时也应记下方向。
\end{enumerate}
在图示中，测得的力用箭头表示，箭头按
声明的比例画出——“每厘米代表十牛顿”——
指向与力相同的方向：这正是童年时的力的箭头，
只是如今有了数量。
\end{method}

\section{重力公式}

\begin{proposition}[重力与质量成正比]\label{prop:g9:weight-vs-mass:formula}
在任一给定地点，物体的重力 $P$（单位为 \unit{N}）
与其质量 $m$（单位为 \unit{kg}）成正比：
\[
P = m \times g ,
\]
其中 $g$ 为该地点的\emph{重力加速度}\index{重力加速度}，
表示每千克物质所受拉力的
牛顿数。在地球表面附近，
\[
g \approx \qty{9.8}{N/kg}:
\]
每千克物质受到的拉力略小于十牛顿。这个
公式是上一章定律的本地形式：$g$ 把地球的
质量与半径合并成一个数，用于表示
“此地”的情况。
\end{proposition}

\begin{omfigure}
\begin{tikzpicture}
\begin{axis}[omaxis, width=10.2cm, height=5.8cm,
    xlabel={质量 $m$（\unit{kg}）}, ylabel={重力 $P$（\unit{N}）},
    xtick={0,1,2,3,4,5}, ytick={0,9.8,19.6,29.4,39.2,49.0},
    yticklabels={$0$,$9.8$,$19.6$,$29.4$,$39.2$,$49.0$},
    ymin=0, ymax=54, xmin=0, xmax=5.5]
  \addplot[very thick, omDef] coordinates {(0,0) (5.2,50.96)};
  \addplot[only marks, mark=*, omThm] coordinates
    {(1,9.8) (2,19.6) (3,29.4) (5,49.0)};
  \node[right, font=\small, omDef] at (axis cs:2.6,20.4)
    {斜率：$g = \qty{9.8}{N/kg}$};
\end{axis}
\end{tikzpicture}

{\small 重力随质量的变化，用测力计与一组已标注的
质量块测得：一条过原点的直线，
其斜率\emph{恰是}该地点的 $g$。}
\end{omfigure}

\begin{example}[公式的应用]\label{ex:g9:weight-vs-mass:atwork}
一名 \qty{60}{kg} 的学生：$P = 60 \times 9.8 = \qty{588}{N}$。一个
\qty{450}{g} 的足球：$P = 0.45 \times 9.8 \approx
\qty{4.4}{N}$——注意换算成千克。反过来，一只箱子
重 \qty{343}{N}，其质量 $m = P/g = 343 \div 9.8 =
\qty{35}{kg}$。再反过来：在一座山上，一个
\qty{10}{kg} 的背包重 \qty{97.7}{N}，则当地的 $g$ 为
$97.7 \div 10 = \qty{9.77}{N/kg}$——$g$ 随海拔
升高而略有减小，正如平方反比规律所预示的。
\end{example}

\begin{example}[同一只箱子，不同的重力]\label{ex:g9:weight-vs-mass:planets}
每个星球都有各自的 $g$，那个古老的宇航员故事
可以写成一张表。你那只忠实的 \qty{10}{kg} 箱子：
\begin{center}
\begin{tabular}{l|c|c}
地点 & $g$（\unit{N/kg}）& 箱子重力 \\
\hline
地球 & $9.8$ & \qty{98}{N} \\
月球 & $1.6$ & \qty{16}{N} \\
火星 & $3.7$ & \qty{37}{N} \\
木星云顶 & $24.8$ & \qty{248}{N}
\end{tabular}
\end{center}
十千克物质到处都一样——但每个星球上的重力各不相同：
质量不变，而重力
随当地的 $g$ 变化。
\end{example}

\begin{remark}[两种秤的最终评判]\label{rem:g9:weight-vs-mass:scales}
童年时的定论，至此以公式作结。天平
比较 $m_1 g$ 与 $m_2 g$：当地的 $g$ 同乘两盘，
这一因子在两盘中相互抵消——天平测得的是\emph{质量}，
在任何星球上都可靠。弹簧测力计直接读出 $P = m g$，
测得的是\emph{重力}，它忠于所在的位置。
浴室秤不过是披着伪装的弹簧测力计，
出厂时以千克标度，方法是用你的体重除以
\emph{地球}的 $g$——在家里诚实可靠，在
月球上却是谄媚者（你每 \qty{60}{kg}，它便显示 \qty{10}{kg}），并且
在木星上则成了诽谤者。
\end{remark}

\section{纸上的力}

\begin{example}[桌上的书，以牛顿计]\label{ex:g9:weight-vs-mass:book}
本书中最古老的平衡，终于用上了数字：一本
\qty{2}{kg} 的字典放在书桌上。它的重力：$P = 2
\times 9.8 = \qty{19.6}{N}$，用一支长 \qty{2}{cm} 的箭头画出，
标度为每厘米十牛顿，竖直向下，指向
书本。桌面的支持力为
\qty{19.6}{N}，竖直向上——一支等长、方向相反的箭头。
这一对力严格相等，静止便是其证明：童年时的
图景，如今成了有度量的账目。
\end{example}

\begin{remark}[力能做什么——缺失的定律]\label{rem:g9:weight-vs-mass:dynamics}
本章测量力，但尚未说明一个
\emph{多余的力}对运动会产生什么作用。童年时的
四种作用——启动、停止、转向、变形——有待其精确定律，
这一定律是物理学中最伟大的定律之一：高中分册
正是以它开启力学。本学年完成的是
打基础的工作：以牛顿为单位的力、由公式算出的重力、按比例绘制的箭头——
这正是即将到来的那句话的语法。
\end{remark}

\section{练习}

\begin{exercise}[$\star$]\label{exo:g9:weight-vs-mass:1}
写出力的单位及其测量工具，并举出两个日常实例，
分别对应一牛顿和一百牛顿的力。
\end{exercise}

\begin{exercise}[$\star$]\label{exo:g9:weight-vs-mass:2}
写出带单位的重力公式，并用文字说明
$g$ 的含义。
\end{exercise}

\begin{exercise}[$\star$]\label{exo:g9:weight-vs-mass:3}
计算地球上的重力：一只 \qty{25}{kg} 的狗；一条 \qty{700}{g}
的面包；以及你自己（自选质量）。
\end{exercise}

\begin{exercise}[$\star$]\label{exo:g9:weight-vs-mass:4}
一个悬挂的瓜使测力计的示数达到 \qty{14.7}{N}，
求它的质量。
\end{exercise}

\begin{exercise}[$\star$]\label{exo:g9:weight-vs-mass:5}
读懂本章的图像：其斜率代表什么，
该直线在月球上会怎样变化？
\end{exercise}

\begin{exercise}[$\star$]\label{exo:g9:weight-vs-mass:6}
一只 \qty{10}{kg} 的行李箱依次游历地球、月球、火星和木星。
表中哪些数值始终不变，哪些会变化，
——原因何在？
\end{exercise}

\begin{exercise}[$\star$]\label{exo:g9:weight-vs-mass:7}
为什么天平在每个星球上都能给出真实的读数，而
浴室秤在月球上却会谄媚？请以公式为依据，而非凭感觉。
\end{exercise}

\begin{exercise}[$\star$]\label{exo:g9:weight-vs-mass:8}
画出（或描述）\qty{2}{kg} 字典所受两个力的箭头，
标度为每厘米五牛顿：写明长度、
方向，并说明这笔力的账是否平衡。
\end{exercise}

\begin{exercise}[$\star\star$]\label{exo:g9:weight-vs-mass:9}
某航空公司规定每件行李限重 \qty{23}{kg}。一位乘客争辩道：
“在月球上它只会重几牛顿！”为什么航空公司的限额——
暂且不论燃料与地板承重——
实际上针对的是质量，应当用什么仪器
在月球机场进行检查？
\end{exercise}

\begin{exercise}[$\star\star$]\label{exo:g9:weight-vs-mass:10}
在火星上，探测车的机械臂重 \qty{111}{N}。它的质量是多少？
同一机械臂回到地球的装配大厅时，
其重力会是多少？
\end{exercise}

\begin{exercise}[$\star\star$]\label{exo:g9:weight-vs-mass:11}
一台测力计平放调零后竖直悬挂使用，示数会略有偏差——
偏向哪一边，为什么？（什么额外的重量会与
所挂物体的重量合在一起？）
\end{exercise}

\begin{exercise}[$\star\star\star$]\label{exo:g9:weight-vs-mass:12}
设计一个实验，画出本章的图像并由此求出 $g$：
写明器材、测量数据表、作图，
以及如何读出斜率。然后说明同一方案在
月球实验室中进行会得到什么结果——
哪些数值会改变，哪些方法依然不受影响。
\end{exercise}

\section{问题：星际装备商}

\begin{problem}[{周末问题 --- 星际装备商的
试衣间；同一物体，四个世界；秤、
测力计与行李台}]\label{pb:g9:weight-vs-mass:1}
装备商为前往太阳系各港口的旅客提供装备
（$g$，单位为 \unit{N/kg}：地球 $9.8$；月球 $1.6$；火星 $3.7$；
木星云层平台 $24.8$）。你负责管理试衣间。

\medskip\noindent\textbf{第一部分 --- 试衣。}
一位质量为 \qty{70}{kg} 的旅客走上前。
\begin{enumerate}
  \item 他们在四个港口各自的重力分别是多少？
  \item 在哪个港口，老式浴室秤
        （刻度为千克，按地球标定）能正确读出其质量？
        ——在其余三个港口，它显示多少？
        （用各港口的重力除以地球的
        $g$。）
  \item 装备商那台诚实的秤是一架
        带标定质量砝码的天平。它在各港口
        分别显示什么？为什么？
  \item 旅客的靴子对地面的压力绝不得
        超过 \qty{2000}{N}（站内地面较薄）。
        有哪个港口，靴子加人体（总质量 \qty{70}{kg}）
        会违反这一规定？
\end{enumerate}

\medskip\noindent\textbf{第二部分 --- 行李台。}
\begin{enumerate}[resume]
  \item 前往火星的行李箱质量为 \qty{40}{kg}，求其
        出发时（地球）与到达时
        （火星）的重力。
  \item 搬运工工会规定，各港口手提的\emph{重力}
        均不得超过 \qty{200}{N}。搬运工在各港口
        最多可携带多大的\emph{质量}？
  \item 一只标有“月球上 \qty{150}{N}”字样的箱子必须
        按质量通过地球海关。请完成换算。
  \item 行李台的测力计显示，悬挂的旅行袋在
        地球上为 \qty{294}{N}。求其质量，以及
        用于货单登记的月球重力。
\end{enumerate}

\medskip\noindent\textbf{第三部分 --- 投诉与
趣闻。}
\begin{enumerate}[resume]
  \item 一位健美运动员从月球归来，怒气冲冲：
        “力量增加六倍，一夜之间全没了！”
        请用公式安慰他：各港口实际改变了什么，
        什么始终未变？
  \item 一位木星平台的厨师抱怨，打蛋时感觉像在举铁砧，
        而打蛋器“质量仅”
        \qty{200}{g}”。求打蛋器在
        木星平台上的重力，并作出裁决。
  \item 礼品店出售“一牛顿巧克力，
        随处可买”。在哪个港口，买家能得到
        最多的巧克力\emph{东西}，并给出
        在地球、月球、木星上各自的质量为多少？
  \item 写出装备商的座右铭：用一句话
        区分随身携带的东西与各个世界
        各个世界收取的费用。
\end{enumerate}
\end{problem}
